\documentclass[11pt,a4paper]{article}
\usepackage{DDTA_DERMATRIAGE_GUIDED_REVIEW_STYLE_R1}

\hypersetup{
  pdftitle={DDTA R25 DermaTriage - Guided Documentation Review R1},
  pdfauthor={Documentation-Driven Threat Analysis research project}
}

\pagestyle{fancy}
\fancyhf{}
\fancyhead[L]{\sffamily\small\color{DDTAGray}DDTA DermaTriage Guided Review}
\fancyhead[R]{\sffamily\small\color{DDTAGray}R25 - WORKING RECORD R1}
\fancyfoot[L]{\sffamily\scriptsize\color{DDTAGray}Documentation reconstruction before Base Analysis}
\fancyfoot[R]{\sffamily\small\color{DDTAGray}\thepage}
\renewcommand{\headrulewidth}{0.4pt}
\renewcommand{\footrulewidth}{0pt}

\begin{document}
\selectlanguage{italian}

% -----------------------------------------------------------------------------
% TITLE
% -----------------------------------------------------------------------------
\begin{titlepage}
\thispagestyle{empty}
\vspace*{8mm}

{\sffamily\bfseries\color{DDTANavy}\fontsize{15}{18}\selectfont Documentation-Driven Threat Analysis}\par
\vspace{5mm}
{\sffamily\bfseries\color{DDTANavy}\fontsize{29}{33}\selectfont DermaTriage Guided Documentation Review}\par
\vspace{2mm}
{\sffamily\color{DDTABlue}\fontsize{18}{22}\selectfont Working Record R1 - dalla source authority alla documentazione DDTA}\par

\vspace{10mm}
\begin{tcolorbox}[enhanced,arc=3pt,boxrule=0pt,colback=DDTANavy,left=12pt,right=12pt,top=11pt,bottom=11pt]
{\color{white}\sffamily\bfseries\large Registro operativo della validazione holdout DermaTriage}\par
\vspace{2mm}
{\color{white!88}\small Applica la guida R4 passo per passo, rende visibile il ragionamento documentale e costruisce progressivamente la baseline da cui, solo dopo il documentation gate, potra' essere derivata la Base Analysis.}
\end{tcolorbox}

\vspace{8mm}
\begin{tabularx}{\textwidth}{L{42mm}Y}
\toprule
\textbf{Stato} & R25 WORKING REVIEW RECORD - DOCUMENTATION PHASE \\
\textbf{Repository} & \texttt{nballestriero/documentation-driven-threat-analysis} \\
\textbf{Baseline verificato} & \texttt{21575a18ccf2feab62b27e3b046fadcd01c90135} \\
\textbf{Metodo} & \texttt{DDTA\_DOCUMENTATION\_BA\_AUTHORING\_GUIDE\_R4} \\
\textbf{Caso holdout} & DermaTriage - ricostruzione da documentazione esistente \\
\textbf{Base Analysis} & BLOCCATA fino al superamento del documentation gate \\
\textbf{Data} & 3 settembre 2026 \\
\bottomrule
\end{tabularx}

\vfill
\begin{ddtaprinciple}[title={Ruolo di questo artifact}]
Questo documento non e' la documentazione finale DermaTriage e non e' una Base Analysis. E' il registro guidato della ricostruzione: conserva source evidence, candidate meaning, gate applicati, esiti, gap e osservazioni metodologiche. Al termine della review dovra' rendere trasparente non soltanto \emph{cosa} e' stato scritto, ma \emph{perche'} ogni elemento DDTA e' stato accettato, modificato, fermato o lasciato non specificato.
\end{ddtaprinciple}
\end{titlepage}

\tableofcontents
\clearpage

% -----------------------------------------------------------------------------
\section{Obiettivo della sessione e disciplina di lavoro}

L'obiettivo concordato e' arrivare alla scrittura della documentazione DermaTriage in modo guidato, usando un progetto reale per rendere osservabile la metodologia DDTA. La review deve inoltre produrre evidence sufficiente per concludere, a valle del lavoro, quali siano i pregi e i difetti di DDTA e in che modo il metodo contribuisca a migliorare la documentazione di progetto.

\begin{ddtacore}[title={Sequenza autorizzata}]
\begin{lstlisting}
repository baseline
    -> methodology reading
    -> original DermaTriage package + SHA-256
    -> authority gate
    -> project problem framing
    -> MacroRequirement
    -> semantic sufficiency
    -> Decision
    -> FunctionalRequirement
    -> Requirement split
    -> Specialized / Security requirements
    -> cross-MR and terminology review
    -> documentation completeness / promotion gate
    -> ONLY THEN Base Analysis
\end{lstlisting}
\end{ddtacore}

La regola collaborativa della sessione e' una sola: nessuna classificazione dell'intero progetto viene eseguita silenziosamente. Ogni elemento viene discusso e chiuso prima di avanzare.

\begin{lstlisting}
source evidence
    -> candidate project meaning
    -> relevant R4 gates
    -> PASS / REWORK / NOT SPECIFIED / CONFLICTING / METHOD PRESSURE
    -> next element only after disposition
\end{lstlisting}

\begin{ddtaanti}[title={Anti-contamination rule}]
Durante questa fase non si usa Base Analysis per decidere project meaning, struttura dei documenti o livello di decomposizione. I vecchi candidate DDTA di DermaTriage possono essere usati in futuro come evidence comparativa, ma non come autorita' sul progetto.
\end{ddtaanti}

% -----------------------------------------------------------------------------
\section{Authority metodologica e repository baseline}

Prima della source review sono stati letti, nell'ordine concordato:
\begin{enumerate}
  \item \texttt{DDTA\_R25\_DOCUMENTATION\_METHOD\_STABILIZATION\_CHECKPOINT\_R1.md};
  \item \texttt{DDTA\_R25\_HOLDOUT\_VALIDATION\_WORK\_PLAN\_R1.md};
  \item \texttt{DDTA\_DOCUMENTATION\_BA\_AUTHORING\_GUIDE\_R4.tex};
  \item \texttt{DDTA\_R25\_DERMATRIAGE\_DOCUMENTATION\_REVIEW\_CONTINUATION\_HANDOFF\_R1.md}.
\end{enumerate}

Il branch \texttt{master} e' stato verificato sul commit:
\begin{lstlisting}
21575a18ccf2feab62b27e3b046fadcd01c90135
\end{lstlisting}

Questo commit e' il continuation baseline della sessione corrente. Il metodo R4 stabilizzato e' il forward authoring protocol; gli esperimenti R5 storici non sono forward authority.

\begin{ddtacheck}[title={Repository authority gate}]
\textbf{PASS.} La sessione parte dal baseline esplicitamente richiesto e non da una ricostruzione cronologica o da un artifact piu' recente scelto per recency.
\end{ddtacheck}

% -----------------------------------------------------------------------------
\section{A0 - Original source package e authority gate}

\subsection{Package pinned}

Il work plan e l'handoff identificano come unica source authority primaria per il significato DermaTriage il package:

\begin{lstlisting}
DermaTriage-Docs-20260830T152637Z-1-001.zip
Expected SHA-256:
E9ED2C507BEFB95F54A52084687CD1E8798863AE81CF69D09568864D8CBF280E
\end{lstlisting}

Il package originale e' stato riaperto e l'hash e' stato ricalcolato sui byte materializzati:

\begin{lstlisting}
Computed SHA-256:
E9ED2C507BEFB95F54A52084687CD1E8798863AE81CF69D09568864D8CBF280E
\end{lstlisting}

\begin{ddtacore}[title={A0 disposition}]
\textbf{PASS.} Nome e SHA-256 coincidono esattamente con il package pinned. Non viene sostituito con archivi simili o con working package prodotti durante precedenti passaggi DDTA.
\end{ddtacore}

\subsection{Source inventory}

L'archivio contiene sette file:

\begin{tabularx}{\textwidth}{L{59mm}Y}
\toprule
\textbf{Source artifact} & \textbf{Ruolo iniziale nella review} \\
\midrule
\nolinkurl{OR2_Architecture_Document.pdf} & Descrizione di purpose, pipeline, integrazioni, API e lifecycle operativo. \\
\nolinkurl{OR2_Model_Test_Report.pdf} & Evidence di modello, metriche, threshold e risultati di test. \\
\nolinkurl{OR3_Dataset_Metadata_Catalog.pdf} & Evidence su dataset, schema e uso dichiarato. \\
\nolinkurl{OR4_Training_Environment_Config.pdf} & Realization/configuration evidence per training e runtime. \\
\nolinkurl{OR4_Training_Cycles_Report.pdf} & Evidence di training e valutazione. \\
\nolinkurl{OR5_Test_Environment_Setup.pdf} & Test/integration evidence, endpoint e pass condition. \\
\nolinkurl{efficientnet_b4.pth} & Binary model evidence; non normative project meaning per default. \\
\bottomrule
\end{tabularx}

\begin{ddtaprinciple}[title={Authority non significa promozione automatica}]
Essere parte del package autoritativo significa che il fatto source-supported deve essere considerato; non significa che ogni fatto tecnico diventi automaticamente \MR, \DEC{} o \REQ. Test, configuration, runtime state e current realization devono essere classificati prima di ricevere status normativo.
\end{ddtaprinciple}

\subsection{A0 gate summary}

\begin{tabularx}{\textwidth}{L{55mm}L{27mm}Y}
\toprule
\textbf{Check} & \textbf{Disposition} & \textbf{Nota} \\
\midrule
Repository baseline & PASS & Commit esatto verificato. \\
Package name & PASS & Nome esatto pinned. \\
Package SHA-256 & PASS & Hash calcolato = hash atteso. \\
Original source as project authority & PASS & Fonte primaria stabilita. \\
Previous DDTA candidates as authority & EXCLUDED & Solo future comparison evidence. \\
Base Analysis as authority & EXCLUDED & Documentation phase ancora aperta. \\
\bottomrule
\end{tabularx}

% -----------------------------------------------------------------------------
\section{A1 - Project problem framing}

\subsection{Source evidence osservata}

La fonte architetturale dichiara esplicitamente:

\begin{ddtaexample}[title={OR2 Architecture Document - purpose}]
\textbf{Purpose:} ``AI-assisted dermatology triage for early skin cancer detection.''
\end{ddtaexample}

La stessa fonte descrive una pipeline che arricchisce il contesto clinico prima della ``final triage decision'', una adaptation layer che assegna una priorita' P1--P4, specialista e SLA, e una successiva doctor validation/correction.

Per il framing, i nomi delle tecnologie e le scelte correnti di soluzione vengono temporaneamente rimossi: EfficientNet, Qwen2-VL, ChromaDB, BioMistral, B4, endpoint, threshold e storage non sono necessari per esprimere il problema generale.

\subsection{Candidate project problem framing}

\begin{ddtacore}[title={Candidate framing accettato nella sessione}]
DermaTriage affronta il problema di supportare il triage precoce di casi dermatologici relativi a lesioni cutanee potenzialmente oncologiche, utilizzando le informazioni disponibili sul caso per determinare una priorita' di presa in carico e supportare l'instradamento verso l'assistenza specialistica appropriata.
\end{ddtacore}

\subsection{Gate D0}

\begin{tabularx}{\textwidth}{L{54mm}L{23mm}Y}
\toprule
\textbf{Domanda} & \textbf{Esito} & \textbf{Motivazione} \\
\midrule
Usa vocabolario di problema/dominio? & PASS & Triage dermatologico, caso, priorita', presa in carico. \\
Resta valido senza la soluzione corrente? & PASS & Nessun modello, provider, API o database e' necessario. \\
Boundary significativo comprensibile? & PASS con gap esplicito & Il triage e' chiaro; l'autorita' diagnostica definitiva non e' stabilita. \\
E' leggibile da un domain reader? & PASS & Non richiede conoscenza dell'architettura corrente. \\
\bottomrule
\end{tabularx}

% -----------------------------------------------------------------------------
\section{Prima semantic-sufficiency question: triage vs diagnosi}

Durante il framing emerge la prima distinzione materiale che il metodo impedisce di nascondere dietro la terminologia della soluzione.

La source documentation usa contemporaneamente:
\begin{itemize}
  \item il purpose di \emph{dermatology triage};
  \item un output Stage 4 chiamato \texttt{pathology};
  \item endpoint e storage denominati \texttt{diagnose}/\texttt{diagnoses};
  \item una successiva doctor validation/correction.
\end{itemize}

\begin{ddtacheck}[title={Smallest critical proposition}]
DermaTriage ha la responsabilita' e l'autorita' di determinare una diagnosi clinica definitiva, oppure la \texttt{predicted\_pathology} e' informazione analitica di supporto al triage e alla successiva review medica?
\end{ddtacheck}

\begin{ddtaopen}[title={Disposition corrente: NOT SPECIFIED}]
La source documentation esaminata non stabilisce in modo sufficiente che DermaTriage abbia autorita' per una diagnosi clinica definitiva. La distinzione non viene completata per inferenza. La \texttt{predicted\_pathology} resta tuttavia un fatto source-supported materiale che dovra' essere conservato e classificato al semantic owner appropriato.
\end{ddtaopen}

\subsection{Diagnosi del problema}

\begin{tabularx}{\textwidth}{L{49mm}L{28mm}Y}
\toprule
\textbf{Candidate diagnosis} & \textbf{Esito} & \textbf{Motivo} \\
\midrule
Source/documentation gap & YES & La fonte usa termini che permettono letture materialmente diverse. \\
Authoring-guide problem & NON DIMOSTRATO & La guida ha fatto emergere la distinzione. \\
Metamodel pressure & NON DIMOSTRATO & Non esiste ancora un significato governato che il metamodel non riesca a rappresentare. \\
BA issue & NON APPLICABILE & BA non e' iniziata. \\
\bottomrule
\end{tabularx}

\begin{ddtaprinciple}[title={Osservazione metodologica 01}]
In questo caso DDTA migliora la documentazione non aggiungendo dettaglio, ma obbligando a separare un termine tecnico presente nell'implementazione da una responsabilita' clinica che la fonte non governa chiaramente. Il valore osservato e' la visibilita' del gap, non la sua chiusura artificiale.
\end{ddtaprinciple}

% -----------------------------------------------------------------------------
\section{A2 - Primo MacroRequirement}

\subsection{Candidate MR-01}

\begin{ddtacore}[title={MR-01 - Valutazione di triage del caso dermatologico}]
\textbf{Intent.} Determinare, dalle informazioni disponibili sul caso dermatologico, una valutazione di urgenza utile a stabilire la priorita' di presa in carico specialistica.

\textbf{Context.} DermaTriage supporta il triage precoce di lesioni cutanee potenzialmente oncologiche.

\textbf{Stakeholders.} Paziente; professionista sanitario coinvolto nella presa in carico e nella review.

\textbf{Scope.} Valutazione del caso fino alla determinazione dell'urgenza/priorita' di triage.

\textbf{Assumptions / Constraints.} Nessuna necessaria al livello MR in questa fase.

\textbf{dependsOn.} Nessuna relazione macro introdotta in questa fase.
\end{ddtacore}

Il candidate MR non incorpora EfficientNet, quattro stage, B4, P1--P4, threshold, pathology prediction o SLA. L'assenza e' intenzionale: questi significati devono essere classificati nei livelli appropriati invece di essere caricati nella macro-responsabilita'.

\subsection{Gate D1 - MR}

\begin{longtable}{L{49mm}L{21mm}L{76mm}}
\toprule
\textbf{Test R4} & \textbf{Esito} & \textbf{Lettura della sessione} \\
\midrule
\endhead
Una sola macro-responsabilita' & PASS & Governa la valutazione di triage del caso. \\
Problem anchoring & PASS & Deriva direttamente dal framing. \\
Architecture / technology resilience & PASS & Sopravvive alla sostituzione completa della soluzione. \\
Temporal stability & PASS & Non dipende dalla configurazione corrente. \\
No comportamento atomico da FR & PASS & Non prescrive stage, chiamate o operazioni. \\
No solution commitment leakage & PASS & Nessun modello, API o provider. \\
Dependency distinta da containment & PASS & Nessuna dependency inventata. \\
Broad domain readability & PASS & Comprensibile senza conoscere DermaTriage internamente. \\
Split image vs symptom-only a livello MR & NO SPLIT & Al momento sono due path della stessa responsabilita' di triage, non due macro-responsabilita' governate. \\
Macro project map globale & PENDING & Potra' essere chiusa solo dopo l'identificazione dell'insieme MR. \\
\bottomrule
\end{longtable}

\subsection{Semantic sufficiency del MR}

La lettura materialmente critica e' ancora la distinzione:

\begin{lstlisting}
triage responsibility
        !=
definitive diagnostic responsibility
\end{lstlisting}

La prima e' \texttt{AFFIRMED}; la seconda resta \NS. Il MR e' scritto in modo da non assorbire la seconda.

Un'ulteriore distinzione source-supported -- urgenza analitica HIGH/MEDIUM/LOW contro priorita' operativa P1--P4 con SLA -- e' materiale, ma non richiede uno split dell'MR. Deve essere conservata nella review downstream.

\begin{ddtacore}[title={MR-01 disposition}]
\textbf{PASS.} Il MR e' semanticamente sufficiente per il proprio livello, con il gap sull'autorita' diagnostica preservato esplicitamente.
\end{ddtacore}

\subsection{STOP test}

\begin{ddtacheck}[title={Serve ulteriore significato governato nello scope corrente?}]
\textbf{YES.} La fonte governa ulteriori distinzioni materiali: path con immagine, fallback symptom-only, pipeline analitica, trasformazione dell'urgenza in priorita' operativa, B4 boundary, clinician review/correction e lifecycle di apprendimento. Questi significati devono essere esaminati uno per volta senza forzare in anticipo il loro owning layer.
\end{ddtacheck}

Pertanto il branch non si ferma a MR-01 e la prossima azione autorizzata e' la review della prima candidate \DEC{} sotto MR-01.

% -----------------------------------------------------------------------------
\section{Stato della review dopo il primo MR}

\begin{center}
\begin{tikzpicture}[
  node distance=7mm and 12mm,
  every node/.style={font=\sffamily\small,align=center},
  done/.style={draw=DDTAGreen,rounded corners=2pt,very thick,fill=DDTALightGreen,minimum width=50mm,minimum height=9mm},
  now/.style={draw=DDTAAmber,rounded corners=2pt,very thick,fill=DDTALightAmber,minimum width=50mm,minimum height=9mm},
  later/.style={draw=DDTAGray,rounded corners=2pt,thick,fill=DDTALightGray,minimum width=50mm,minimum height=9mm},
  blocked/.style={draw=DDTARed,rounded corners=2pt,very thick,fill=DDTALightRed,minimum width=50mm,minimum height=9mm},
  arr/.style={-{Latex[length=2.2mm]},thick,draw=DDTANavy}
]
\node[done] (a0) {A0 Authority / source gate\\PASS};
\node[done,below=of a0] (a1) {A1 Problem framing\\PASS + explicit NOT SPECIFIED};
\node[done,below=of a1] (a2) {A2 MR-01 + semantic sufficiency\\PASS};
\node[now,below=of a2] (a4) {NEXT: first Decision candidate\\gate-by-gate review};
\node[later,below=of a4] (rest) {FR / split / SR / SecR / boundaries\\terminology / propagation / completeness};
\node[blocked,below=of rest] (ba) {Base Analysis\\BLOCKED until documentation gate};
\draw[arr] (a0)--(a1);
\draw[arr] (a1)--(a2);
\draw[arr] (a2)--(a4);
\draw[arr] (a4)--(rest);
\draw[arr] (rest)--(ba);
\end{tikzpicture}
\end{center}

\begin{tabularx}{\textwidth}{L{42mm}L{28mm}Y}
\toprule
\textbf{Review object} & \textbf{Status} & \textbf{Disposition} \\
\midrule
Repository baseline & CLOSED & Exact commit verified. \\
Original package & CLOSED & Exact SHA-256 verified. \\
Authority gate & CLOSED & Original source is project authority. \\
Problem framing & CLOSED for current scope & Triage framing accepted; diagnostic authority remains \NS. \\
MR-01 & CLOSED for current scope & PASS; branch continues. \\
First Decision & NEXT & Not yet authored. \\
Documentation gate & OPEN & Most conditions not yet evaluated. \\
Base Analysis & BLOCKED & Explicitly not started. \\
\bottomrule
\end{tabularx}

% -----------------------------------------------------------------------------
\section{Protocollo guidato per i prossimi elementi}

Questa sezione definisce il formato con cui il documento dovra' essere esteso durante la continuazione. L'obiettivo e' che ogni decisione documentale resti auditabile.

\subsection{Worksheet per ogni candidate DDTA element}

\begin{ddtacheck}[title={1. Source evidence}]
Quali frasi, tabelle, diagrammi o altri artifact del package sostengono il significato? L'evidence descrive project responsibility/commitment oppure soltanto current realization, configuration, runtime o test state?
\end{ddtacheck}

\begin{ddtacheck}[title={2. Candidate project meaning}]
Qual e' la minima proposizione governata che vogliamo preservare? E' scritta con vocabolario di progetto e al livello di astrazione pertinente?
\end{ddtacheck}

\begin{ddtacheck}[title={3. Semantic owner}]
Il significato appartiene a framing, \MR, \DEC, \FR, \SR, \SECR, boundary cross-MR, oppure e' supporting realization/evidence senza collocazione normativa automatica?
\end{ddtacheck}

\begin{ddtacheck}[title={4. Relevant gates}]
Applicare tutti i gate pertinenti, non soltanto quello del tipo candidato. Per una Decision, ad esempio, includere responsibility-boundary e Decision-vs-implementation evidence. Per un FR includere parentage, coherent-unit/split, binding e failure semantics quando materiali.
\end{ddtacheck}

\begin{ddtacheck}[title={5. Explicit disposition}]
Registrare una delle disposition operative: \texttt{PASS}, \texttt{REWORK}, \texttt{NOT SPECIFIED}, \texttt{CONFLICTING}, \texttt{OUTSIDE CURRENT DDTA DOCUMENTATION}, \texttt{METHOD PRESSURE}.
\end{ddtacheck}

\begin{ddtacheck}[title={6. Stopping decision}]
Se il significato corrente e' sufficiente, fermare il branch. Non creare elementi inferiori per completare visivamente una gerarchia.
\end{ddtacheck}

\subsection{Ordine rimanente della review}

\begin{longtable}{L{17mm}L{39mm}L{90mm}}
\toprule
\textbf{Step} & \textbf{Review} & \textbf{Domanda guida} \\
\midrule
\endhead
A3 & Semantic sufficiency & Sopravvivono letture materialmente diverse? Qual e' la smallest critical proposition? \\
A4 & Decision & Quale commitment materiale restringe l'MR? E' governato o solo implementation evidence? \\
A5 & FunctionalRequirement & Quale obbligo operativo independently assessable discende dalla Decision? \\
A6 & Requirement split & Una clausola puo' cambiare o essere assessed indipendentemente dalle altre? \\
A7 & SpecializedRequirement & Esiste una proprieta' concern-specific aggiuntiva che strengthening il parent FR? \\
A8 & SecurityRequirement & Esiste una singola security property governata con failure mode identificabile? \\
A9 & Cross-MR / consumed service & Il branch possiede o consuma la capability? Quale garanzia del servizio e' davvero governata? \\
A10 & Terminology / bindings & Le stesse identita' usano termini e riferimenti stabili? \\
A11 & Propagation & Le correzioni upstream sono preservate nei discendenti senza narrowing/widening silenzioso? \\
A12 & Completeness / promotion & La candidate baseline e' completa, coerente e accettabile come source baseline per BA? \\
\bottomrule
\end{longtable}

% -----------------------------------------------------------------------------
\section{Registro delle osservazioni metodologiche}

Le osservazioni in questa sezione sono volutamente separate dalla documentazione di progetto. Non devono contaminare MR/Decision/Requirement finali. Verranno consolidate solo al termine della documentazione e, successivamente, confrontate con la regressione Documentation $\leftrightarrow$ BA.

\subsection{Pregi osservati finora - provvisori}

\begin{enumerate}
  \item \textbf{Authority discipline.} La metodologia impedisce di prendere come verita' il candidate DDTA precedente, il modello corrente o la BA futura. Questo rende piu' chiaro da dove proviene ogni significato.
  \item \textbf{Problem/solution separation.} Il framing costringe a rimuovere tecnologia e architettura quando non sono intrinseche al problema. Il risultato e' una responsabilita' piu' stabile nel tempo.
  \item \textbf{Ambiguity surfacing.} La distinzione triage vs diagnosi definitiva, facilmente nascosta da endpoint e campi chiamati \texttt{diagnose}/\texttt{pathology}, viene resa esplicita come \NS{} invece di essere risolta per supposizione.
  \item \textbf{Controlled decomposition.} Il gate di STOP impedisce di produrre documenti inferiori soltanto per riempire la gerarchia, mentre il parentage mantiene la tracciabilita' quando la decomposizione e' necessaria.
  \item \textbf{Classification before method change.} Un problema incontrato viene prima classificato come source gap, guide problem, metamodel pressure, L2 issue o downstream issue. Questo riduce il rischio di cambiare il metodo per ogni anomalia documentale.
\end{enumerate}

\subsection{Costi, limiti e pressure da verificare - non conclusivi}

\begin{enumerate}
  \item \textbf{Review effort.} L'applicazione esplicita dei gate richiede piu' lavoro di una riscrittura informale della documentazione. Dovremo valutare se il maggior costo produce un beneficio proporzionato in chiarezza e regressione.
  \item \textbf{Judgement burden.} La distinzione fra project commitment e current realization non e' completamente automatizzabile; il metodo richiede judgement governato e quindi disciplina del reviewer.
  \item \textbf{Technical-realization preservation.} Dataset, modelli, threshold, training configuration e test evidence possono essere materialmente utili senza essere Requirements o Decisions. La loro rappresentazione non normativa resta una pressure area aperta nella R4.
  \item \textbf{Risk of procedural verbosity.} Se il registro dei gate viene confuso con la documentazione finale, il metodo puo' produrre troppo research commentary. La separazione fra working review record e project baseline sara' quindi un test importante.
  \item \textbf{Tooling opportunity with authority boundary.} Molti check strutturali possono essere assistiti da tool, ma semantic owner, equivalenza e meaning non devono essere decisi automaticamente senza governed review.
\end{enumerate}

\subsection{Come DDTA sta gia' migliorando la documentazione DermaTriage}

\begin{tabularx}{\textwidth}{L{52mm}Y}
\toprule
\textbf{Problema tipico nella fonte classica} & \textbf{Miglioramento prodotto dalla review DDTA} \\
\midrule
Purpose, architecture e implementation facts sono mescolati. & Il framing separa il problema stabile dalle scelte correnti di soluzione. \\
Terminologia come \texttt{diagnose}, \texttt{pathology} e ``triage'' puo' suggerire authority diverse. & La semantic sufficiency isola la proposizione critica e conserva \NS{} invece di scegliere una lettura. \\
Una pipeline tecnica suggerisce naturalmente una decomposizione per componenti. & L'MR viene definito per responsibility, non per stage architetturale. \\
Metriche, threshold e configuration sembrano ``importanti'' e quindi rischiano di diventare requisiti per default. & La parameter/evidence classification richiede prima di stabilirne il semantic owner. \\
La mancanza di informazione puo' essere nascosta da wording generico. & \NS{} e \texttt{CONFLICTING} diventano esiti espliciti e tracciabili. \\
\bottomrule
\end{tabularx}

\begin{ddtaworking}[title={Criterio per la valutazione finale della metodologia}]
Alla fine del caso DermaTriage non giudicheremo DDTA in base a quanto ``bella'' appare la gerarchia. Valuteremo invece se il metodo ha: preservato il significato source-supported; reso visibili gap e conflitti; ridotto ambiguity e implementation leakage; prodotto una baseline comprensibile da un reviewer di progetto; consentito una BA fedele senza invenzione; richiesto un costo di authoring/review ragionevole rispetto ai benefici.
\end{ddtaworking}

% -----------------------------------------------------------------------------
\section{Forma prevista della documentazione finale DermaTriage}

Questo working record non impone in anticipo quali documenti debbano esistere. La forma finale sara' determinata dai gate. Tuttavia, se i branch richiederanno l'intera decomposizione corrente, la relazione di ownership restera':

\begin{lstlisting}
Project problem framing
    -> MacroRequirement
        -> Decision
            -> FunctionalRequirement
                -> SpecializedRequirement
                    -> SecurityRequirement
\end{lstlisting}

Ogni branch potra' fermarsi prima quando semanticamente sufficiente. La documentazione finale dovra' essere piu' corta e piu' leggibile di questo registro: conterra' project meaning governato, non il processo di ricerca che lo ha prodotto.

\begin{ddtaprinciple}[title={Working record vs project baseline}]
\textbf{Working record:} evidence, dubbi, gate, alternative, diagnostics, method observations.\par
\textbf{Project baseline finale:} solo significato governato, gap espliciti necessari, termini stabili, responsabilita', commitment e obblighi accettati.
\end{ddtaprinciple}

% -----------------------------------------------------------------------------
\section{Prossimo passo autorizzato}

La sessione e' arrivata esattamente al punto seguente:

\begin{ddtaworking}[title={NEXT - first Decision under MR-01}]
Prendere una sola candidate Decision sotto MR-01 e applicare insieme:
\begin{enumerate}
  \item source evidence;
  \item candidate commitment;
  \item Decision gate;
  \item responsibility-boundary gate;
  \item Decision-vs-current-realization test;
  \item SELECTABLE / NECESSITY-CONSTRAINED / DEFAULT review;
  \item semantic-sufficiency check;
  \item explicit disposition e STOP decision.
\end{enumerate}
Nessun FR viene creato prima che questa review sia chiusa.
\end{ddtaworking}

% -----------------------------------------------------------------------------
\appendix
\section{Compact review worksheet}

\begin{longtable}{L{45mm}L{105mm}}
\toprule
\textbf{Field} & \textbf{Da compilare per ogni elemento} \\
\midrule
\endhead
Source artifact / location & Documento, sezione/pagina, tabella o artifact che sostiene il significato. \\
Source proposition & Fatto minimo effettivamente supportato. \\
Candidate DDTA type & MR / Decision / FR / SR / SecR / boundary / evidence-only / unresolved. \\
Candidate wording & Testo minimo nel vocabolario del progetto. \\
Semantic owner & Perche' questo livello e non un altro. \\
Relevant gates & Tutti i gate applicati, non solo quello nominale del tipo. \\
Smallest critical proposition & Distinzione minima quando esistono letture materialmente diverse. \\
Evidence status & AFFIRMED / DENIED / NOT SPECIFIED / CONFLICTING. \\
Disposition & PASS / REWORK / NOT SPECIFIED / CONFLICTING / OUTSIDE / METHOD PRESSURE. \\
STOP decision & Fermare il branch o continuare, con motivo. \\
Method observation & Eventuale OBS-DDTA / HYP-METHOD / MIP / DT-METHOD candidate, separata dal project meaning. \\
\bottomrule
\end{longtable}

\section*{Session snapshot}

\begin{lstlisting}
BASELINE
21575a18ccf2feab62b27e3b046fadcd01c90135

SOURCE PACKAGE
DermaTriage-Docs-20260830T152637Z-1-001.zip
SHA-256
E9ED2C507BEFB95F54A52084687CD1E8798863AE81CF69D09568864D8CBF280E

A0 AUTHORITY GATE                  PASS
A1 PROJECT PROBLEM FRAMING         PASS
DIAGNOSTIC AUTHORITY               NOT SPECIFIED
A2 MR-01                           PASS
MR-01 STOP TEST                    CONTINUE
FIRST DECISION                     NEXT
DOCUMENTATION GATE                 OPEN
BASE ANALYSIS                      NOT STARTED / BLOCKED
\end{lstlisting}

\begin{ddtacore}[title={Continuation invariant}]
La prossima revisione di questo documento deve aggiungere evidence e gate senza riscrivere retroattivamente gli esiti gia' accettati, salvo esplicita re-open motivata da nuova source evidence o da una contraddizione rilevata durante la regression.
\end{ddtacore}

\end{document}
